\documentclass[17pt]{extarticle}
\usepackage{graphicx}
\usepackage{amsmath}
\usepackage{amssymb}
\usepackage{hyperref}
\title{Math Expression}
\author{Zadid Bin Azad}
\date{\today}

\begin{document}

\maketitle
\newpage

\section{Inline Math Mode}

The quadratic formula \(x= \frac{-b \pm \sqrt{b^2-4ac}}{2a}\) was invented way back in the Egyptian time.
\\
In \LaTeX shorthand there is another way to write inline equations that is $E = mc^2$
\\
For the third method, let's say area of a circle is \begin{math}
Area = \pi R^2
\end{math}

\section{Display Math Mode}
This is the un-numbered mode.
\[\int_0^\infty e^{-x^2}\,dx = \frac{\sqrt{\pi}}{2}\]
.
The \LaTeX shorthand method is next; 
$$
f(x) = \sum_{n=0}^{\infty} \frac{x^n}{n!}
$$

\begin{displaymath}
\lim_{x \to 0} \frac{\sin x}{x} = 1
\end{displaymath}

\section{Equation Environment}
There are two variants, one is equation and equation*

\begin{equation} \label{eqn:pythagoras}
    a^2 + b^2 = c^2
\end{equation}

We can refer to an equation with \autoref{eqn:pythagoras}.

Same equations without numbering
\begin{equation*}
    a^2 + b^2 = c^2
\end{equation*}


\section{Math functions, Superscripts \& Subscripts}
$x^2,\quad x^{10}, \quad x_i, \quad x_{ij}, \quad \sum_{i=1}^{n} x_i$ 
\\
Trigonometric Function:
$\sin{x^2}, \quad \cos{\theta},\quad \tan^{-1}{x}$
\\
nth Root:
$\sqrt[3]{8}$
\\
\\
$\frac{a}{b} \times \frac{c}{d}$

\section{Symbols}
$$\neq$$
$$\leq$$
$$\geq$$
$$\approx$$
$$\equiv$$
$$\propto$$
$$\ll$$
$\gg$

\section{Brackets}
\[
    \left( \frac{\frac{ac}{d^2}}{\sin{\frac{2\theta}{\pi}}} \right)
\]


\section{Greek letters}
$$\alpha$$
$$\Delta$$
$$\delta$$
$$\phi$$
$$\Phi$$
$$\mu$$

Special Math fonts
$$\mathcal{A}$$
$$\mathfrak{H}$$
$$\mathbb{R}$$
$$\mathbb{Z}$$
$$\mathbb{N}$$

\section{Vectors}
$$\Vec{v}$$
$$\Vec{AB}$$

$$\vec{u} \cdot \vec{v} = |\vec{u}||\vec{v}|\cos{\theta}$$

\section{Aligning Equations}
\begin{align}
    x + y &= 5 \\
    2x - y &= 10 \\
    x &=2
\end{align}

\begin{align*}
    (a+b)^2 &= a^2 + 2ab + b^2 \\
    &= a^2 +b^2 + 2ab
\end{align*}

\begin{gather*}
    a+b = c \\
    x^2 + y^2 =r^2
\end{gather*}

\begin{multline}
    p(x) = x^8 + x^7 + x^6 + x^8 + x^7 + x^6  \\
    +x^5 + x^4 + x^8 + x^7 + x^6 +1
\end{multline}

\begin{equation}
    \begin{split}
        f(x) &= x^2 + 2x+1\\
        &= (x+1) ^2
    \end{split}
\end{equation}


\section{Matrices}
\[
\begin{matrix}
    a & b \\
    c & d
\end{matrix}
\]

\[
\begin{pmatrix}
    a & b \\
    c & d
\end{pmatrix}
\]

\[
\begin{bmatrix}
    a & b \\
    c & d
\end{bmatrix}
\]

\[
\begin{Bmatrix}
    a & b \\
    c & d
\end{Bmatrix}
\]

\[
\begin{vmatrix}
    a & b \\
    c & d
\end{vmatrix}
\]

\[
\begin{Vmatrix}
    a & b \\
    c & d
\end{Vmatrix}
\]

\end{document}